\subsection{单代数的单子代数}

定理 5. —— 设 A 是一个中心单 K-代数，B 是 A 的一个在 K 上为有限次数的半单子代数。

\begin{enumerate}
\def\labelenumi{\alph{enumi})}
\item
  B 在 A 中的交换子 B' 是一个单子代数，且 B 是 B' 在 A 中的交换子。此外，代数 \(B \cap B'\) 是一个在 K 上为有限次数的半单交换代数；它是 B 与 B' 的公共中心。
\item
  假设 B 是单的。则 \(B'\) 是单的，并且我们有等式
\end{enumerate}

\[
[ \mathrm{A}: \mathrm{B} ^ {\prime} ] _ {s} = [ \mathrm{B}: \mathrm{K} ], \quad [ \mathrm{A}: \mathrm{B} ] _ {s} = [ \mathrm{B} ^ {\prime}: \mathrm{K} ], \quad [ \mathrm{A}: \mathrm{K} ] = [ \mathrm{B}: \mathrm{K} ] [ \mathrm{B} ^ {\prime}: \mathrm{K} ].
\]

（左次数 \([A : B]_s\) 的定义参见 VIII, 第 124 页，定义 2。）

K-代数 \(A^{\circ}\) 是中心单的，而 K-代数 B 是半单的且为有限次数的。由 VIII, 第 221 页的推论 1，代数 \(C = B \otimes_{K} A^{\circ}\) 是半单的。设 M 是这样的 C-模：它与 A 有相同的加群，其外部律由公式 \((b \otimes a)a' = ba'a\) 给出，其中 \(a, a'\) 属于 A 而 b 属于 B。设 u 是 \(\operatorname{End}_{\mathbf{Z}}(\mathrm{A})\) 的一个元素。则 u 属于 C-模 M 的交换子 \(C_{M}'\) 当且仅当 u 是右 A-线性的且是左 B-线性的，换言之，当且仅当 u 属于 \(B_{M}\) 在 A-模 \(A_{s}\) 的位似环中的交换子。因此，我们定义一个同构 \(\gamma\)（从 \(B'\) 到 \(C_{M}'\) 的），它由关系式 \(\gamma(b')(x) = b'x\)（其中 \(b'\) 属于 \(B'\)，x 属于 M）给出。现在，环 C 是半单的，且 C-模 M 由 A 的元素 1 生成。由 VIII, 第 139 页的命题 6，环 \(C_{M}'\) 是半单的，于是代数 \(B'\) 是半单的。

设 \(\varphi\) 是从 \(A \otimes_K A^\circ\) 到 \(\mathrm{End}_K(M)\) 的 K-代数同态，它把 \(a \otimes a'\) 映为从 M 到 M 的 K-线性映射 \(x \mapsto axa'\)。由于 K-代数 A 与 \(A^\circ\) 都是中心单的，\(A \otimes_K A^\circ\) 的仅有的双边理想是 0 与 \(A \otimes_K A^\circ\)（VIII, 第 218 页，推论）。我们有 \(C_M = \varphi(B \otimes A^\circ)\) 与 \(C_M' = \varphi(B' \otimes K)\)。同态 \(\varphi\) 非零；因而它是单射。由于环 C 是半单的，由 VIII, 第 139 页的命题 5，我们有 \(C_M'' = C_M\)。由此推出，子代数 \(B \otimes_K A^\circ\)（\(A \otimes_K A^\circ\) 的）是子代数 \(B' \otimes_K K\) 的交换子。于是 \(B' \otimes_K K\) 在 \(A \otimes_K K\) 中的交换子等于 \((B \otimes_K A^\circ) \cap (A \otimes_K K)\)，即由 II, §7, No.~9, 第 311 页的命题 19，等于 \(B \otimes K\)。因此 \(B'\) 在 A 中的交换子等于 B。代数 \(L = B \cap B'\) 是 B 的中心。由于 B 是 K 上一个有限次数的半单代数，代数 L 是交换的、半单的，且在 K 上是有限次数的。由于 B 是 \(B'\) 在 A 中的交换子，\(B'\) 的中心也等于 \(L = B \cap B'\)（VIII, 第 77 页）。我们证明了 a)。

现在假设代数 B 是单的。由 VIII, 第 222 页的推论 2，环 C 是单的。把 VIII, 第 123 页的命题 4 应用于其交换子同构于 \(B'\) 的 C-模 M，便得环 \(B'\) 是单的且 M 是一个有限长度的 \(B'\)-模。换言之，\(B'\) 是单环 A 的一个单子环，且左次数 \([A : B']_{s}\) 是一个整数 \(m \geqslant 1\)。作为左 \(B'\)-模，A 有一个有限基 \((a_{1}, \ldots, a_{m})\)。此外（出处同前），\(\varphi\) 通过限制诱导出从 C 到 \(\mathrm{C}_{\mathrm{M}}'' = \mathrm{End}_{\mathrm{B}'}(\mathrm{A})\) 的一个同构，于是映射 \(c \mapsto (ca_{1}, \ldots, ca_{m})\)（从 C 到 \(A^{m}\) 的）是双射。因此，C 是一个维数为 m 的自由右 A-模。现在，我们有 \(C = B \otimes A^{o}\)，故 C 是一个维数为 {[}B : K{]} 的自由右 A-模；由此推出 \([A : B']_{s} = m = [B : K]\)。由 VIII, 第 125 页的命题 6，我们导出

\[
[ \mathrm{A}: \mathrm{K} ] = [ \mathrm{A}: \mathrm{B} ^ {\prime} ] _ {s} [ \mathrm{B} ^ {\prime}: \mathrm{K} ] = [ \mathrm{B}: \mathrm{K} ] [ \mathrm{B} ^ {\prime}: \mathrm{K} ];
\]

由于我们还有

\[
[ \mathrm{A}: \mathrm{K} ] = [ \mathrm{A}: \mathrm{B} ] _ {s} [ \mathrm{B}: \mathrm{K} ]
\]

且 \([B:K]\) 有限非零，我们得出等式 \([A:B]_{s}=[B^{\prime}:K]\)（集合论, III, §6, No.~3, 第 188 页，推论 3）。我们证明了 b)。

\pandocbounded{\includegraphics[keepaspectratio]{images/part002_69e19a74.jpg}}

设 A 是一个有限次数的中心单 K-代数。可以存在 A 的半单交换子代数 B 满足 \([A:K]\neq[B:K][B':K]\)（VIII, 第 269 页，习题 1）。

定理 6. —— 设 A 是一个中心单 K-代数，B 是 A 的一个有限次数的子代数，B' 是它在 A 中的交换子。

\begin{enumerate}
\def\labelenumi{\alph{enumi})}
\item
  假设 B 是中心单的。则 \(B'\) 是中心单的，且 K-代数同态 \(\theta : B \otimes_{K} B' \to A\)（它把 \(b \otimes b'\) 映为 \(bb'\)）是一个同构。
\item
  假设 B 是半单的，并设 \(L = B \cap B'\)。则 \(B'\) 是一个半单代数。L 在 A 中的交换子 \(L'\) 是一个以 L 为中心的半单环，且环同态 \(\psi : B \otimes_{L} B' \to L'\)（它把 \(b \otimes b'\) 映为 \(bb'\)）是一个同构。
\end{enumerate}

我们来证明 a)。若 B 是中心单的，则由 VIII, 第 259 页的定理 5，\(B'\) 是中心单的。于是 K-代数 \(B \otimes_{K} B'\) 是单的（VIII, 第 222 页，推论 2），且同态 \(\theta : B \otimes_{K} B' \to A\) 是单射。现在，由 \([A : B']_{s}\) 与 \([B : K]\) 的相等（定理 5），左 \(B'\)-模 \(B \otimes_{K} B'\) 与 A 是具有相同有限维数的自由模；因而它们是具有相同有限长度的 \(B'\)-模。由 II, §1, No.~10, 第 213 页的推论 2，\(\theta\) 是双射。

我们来证明 b)。由定理 5，代数 L 是交换的、在 K 上是有限次数的且是半单的。把同上应用于 L，其在 A 中的交换子 \(L'\) 是一个半单代数，且 L 是 \(L'\) 在 A 中的交换子，故 L 是 \(L'\) 的中心。由于 L 是半单环 \(L'\)、B 与 \(B'\) 的中心，我们可以把 \(L'\) 等同为单环 \(L_{i}'\)（\(i \in I\)）的一个有限积，使得我们有

\[
\mathrm{L} = \prod_ {i \in \mathrm{I}} \mathrm{L} _ {i}, \qquad \mathrm{B} = \prod_ {i \in \mathrm{I}} \mathrm{B} _ {i}, \qquad \mathrm{B} ^ {\prime} = \prod_ {i \in \mathrm{I}} \mathrm{B} _ {i} ^ {\prime},
\]

其中 \(L_{i}\) 是 \(L_{i}^{\prime}\) 的中心，而 \(B_{i}\) 与 \(B_{i}^{\prime}\) 是 \(L_{i}^{\prime}\) 的以 \(L_{i}\) 为中心的子代数，它们在 \(L_{i}^{\prime}\) 中互为交换子。我们把 \(L_{i}^{\prime}\) 视为交换域 \(L_{i}\) 上的一个中心单代数，把 \(B_{i}\) 视为 \(L_{i}\) 上的一个有限次数的中心单代数。由断言 a)，典范映射 \(\psi_{i}: B_{i} \otimes_{L_{i}} B_{i}^{\prime} \to L_{i}^{\prime}\)（它把 \(b_{i} \otimes b_{i}^{\prime}\) 映为 \(b_{i}b_{i}^{\prime}\)）是一个环同构。现在，我们可以把 \(B \otimes_{L} B^{\prime}\) 等同为 \(\prod_{i \in I}(B_{i} \otimes_{L_{i}} B_{i}^{\prime})\)，于是 \(\psi\) 是映射族 \((\psi_{i})_{i \in I}\) 的积。因此，\(\psi\) 是一个环同构。

推论. —— 假设域 K 是代数封闭的，且 A 是 K 上一个有限次数的单代数。设 B 是 A 的一个单子代数，B' 是 B 在 A 中的交换子。则 B' 是一个单 K-代数，B 是 B' 的交换子，我们有 \([A:K]=[B:K][B':K]\)，且从 \(B\otimes_{K}B'\) 到 A 的典范同态是一个 K-代数同构。

由于 \(K\) 上每个有限次数的单代数都是中心的，该推论由定理 5 与定理 6 得出。

